%%
% BIThesis 研究生学位论文模板 The BIThesis Template for Graduate Thesis
% This file has no copyright assigned and is placed in the Public Domain.
%%

\chapter{绪论}

\label{chap:intro}

\section{研究背景与意义}

矿井巷道环境是典型的卫星导航拒止环境，车辆、巡检机器人以及无人化运输装备的位姿获取主要依赖车载传感器自主完成。随着智慧矿山和矿山无人化技术的持续推进，矿井场景中的高精度环境建图、稳定定位以及大尺度先验地图服务能力，已成为保障无人装备安全运行和提高生产效率的重要基础。特别是在长距离运输、巡检监测、灾害预警和应急救援等任务中，若定位结果出现明显漂移或地图调用效率不足，将直接影响路径跟踪、作业调度和系统鲁棒性。

与地面开放场景相比，矿井巷道具有空间狭长、结构重复、光照条件差、纹理稀缺以及分支拓扑复杂等特点。一方面，三维激光雷达能够提供较稳定的几何观测信息，因而成为矿井场景定位与建图的重要传感器；另一方面，巷道环境沿延伸方向通常存在较强的结构相似性，而在横断面方向又呈现相对稳定的边界约束。这种场景特性使得仅依赖局部点云配准的建图与定位方法容易在长距离运行中产生误差累积，在结构重复区域还可能出现匹配退化和全局一致性不足等问题。

围绕矿井巷道场景的定位任务，当前研究通常涉及三个相互关联的层面。其一，离线建图阶段需要获得全局一致且结构准确的高精先验地图，为后续定位提供可靠的参考基础。其二，在线定位阶段需要在先验地图约束下融合局部运动信息与全局地图观测，在结构重复和观测质量波动条件下保持稳定的位姿估计。其三，当先验地图范围不断增大时，还需要进一步解决整图加载代价高、地图调用效率低以及持续运行资源压力大的问题。若这三个层面无法有效衔接，则高精地图难以真正服务于在线定位任务。

因此，针对矿井巷道场景开展“高精地图构建--先验约束定位--地图组织调用”一体化研究，既具有明确的理论意义，也具有较强的工程应用价值。从理论层面看，挖掘矿井巷道中的几何结构规律，并将其用于建图与定位建模，有助于扩展复杂地下环境中多源结构先验的利用方式；从应用层面看，构建适用于矿井场景的高精地图支撑定位方法，有助于提升无人矿车和井下机器人在真实作业环境中的自主运行能力。

\section{国内外研究现状}

\subsection{三维激光建图与位姿图优化研究现状}

三维激光建图是移动机器人环境感知中的基础问题，其核心在于利用连续点云观测估计载体运动，并在统一坐标系下构建全局地图。现有研究通常采用“前端配准+后端优化”的基本框架：前端通过相邻帧或局部子地图之间的点云配准获得相对运动观测，后端则利用回环约束、先验约束或图优化方法抑制误差累积并提高全局一致性。围绕该框架，国内外学者已在回环检测、平面结构建模、稀疏特征匹配以及鲁棒配准等方面开展了大量研究\cite{Hess2016LoopClosure,Ma2016CPASLAM,Rusu2009FPFH,Yang2021TEASER}。

对于一般城市场景或室内场景，基于回环检测和图优化的建图方法通常能够取得较好的效果。然而，矿井巷道环境在长距离运行中往往缺少稳定且高频的回环机会，单纯依赖局部点云配准或稀疏回环约束，难以从全局尺度有效抑制轨迹漂移。尤其是在长直巷道、转弯连接区域和结构重复区域，局部观测的几何约束强度并不均衡，若缺少对巷道长度关系、方向关系以及局部平面结构的进一步利用，容易造成地图全局一致性不足。

近年来，面向结构化场景的建图研究开始关注显式几何先验的引入，例如平面、线段、方向约束和语义结构约束等。这类方法说明，仅依赖局部点云匹配难以充分发挥场景中的稳定结构信息，而将可重复利用的全局几何关系显式建模到优化过程中，有助于提高建图质量。对于矿井巷道场景而言，如何从环境自身的几何规律出发，构建适用于长距离巷道建图的结构约束模型，仍然是值得进一步研究的问题。

\subsection{先验地图约束定位研究现状}

在已有高精先验地图条件下，定位问题通常转化为当前观测与先验地图之间的配准与状态估计问题。现有研究普遍采用“运动预测+局部配准+地图配准”的技术路线：首先利用 IMU、里程计或历史位姿生成初始预测，再通过帧间点云配准获得局部连续运动约束，最后结合先验地图进行全局匹配，以获得稳定的位姿估计结果。围绕点云配准模型、观测权重分配与多源信息融合，已经形成了较为丰富的研究基础。

然而，矿井巷道环境中的先验地图定位与开放场景存在明显差异。由于巷道沿延伸方向的结构重复性较强，不同位置之间可能出现局部几何相似而全局位置不同的情况；同时，巷道横断面方向通常存在较稳定的边界结构，使得不同方向上的约束强度呈现出不均衡特征。若仍采用统一置信度或统一权重处理帧间配准观测与地图配准观测，则容易在局部观测退化、地图匹配得分波动或初值偏差增大时引起定位不稳定。

为提高定位系统在复杂环境中的鲁棒性，近年来研究者开始关注先验地图约束、多源观测一致性判别以及差异化融合等问题。相关工作表明，仅将先验地图作为被动的匹配对象往往不足以充分发挥其结构约束能力，若能够进一步利用地图中的结构分布特征、局部区域可靠性以及方向性信息，则有望增强定位结果在复杂场景中的稳定性。对于矿井巷道场景而言，如何利用先验地图中的空间占据信息与巷道方向结构，对不同来源观测实施差异化约束，仍有进一步深入研究的必要。

\subsection{大尺度先验地图组织与调用研究现状}

随着定位范围的扩大，高精先验地图的组织方式逐渐成为影响系统在线性能的重要因素。在大尺度环境中，若将整幅地图长期常驻内存并在每一时刻参与匹配，不仅会带来较高的内存占用，还会显著增加地图读取、点云拼接与匹配计算的开销。因此，在自动驾驶、移动机器人和大规模环境感知等领域，地图分块、局部子图调度、瓦片化管理以及按需加载等方法得到了广泛关注。

从现有研究看，地图组织与调用方法通常包含两个关键问题：一是如何将全局地图转化为便于在线检索与组合的结构化单元；二是如何依据当前位置和运行状态稳定地选择局部地图范围，避免边界附近的频繁切换。对于规则道路场景或室外开放场景，这类方法已表现出较好的实用价值。但矿井巷道环境具有狭长、多分支、连续延伸等特征，其地图调用往往与当前位置、巷道拓扑以及局部观测覆盖范围更紧密相关，若缺乏适配矿井场景特点的组织与调用机制，则难以兼顾地图使用效率和定位稳定性。

综合来看，现有研究已分别在建图、定位和地图服务等方面形成了一定基础，但对于矿井巷道这一典型地下结构化场景，仍存在以下不足：其一，离线建图对巷道全局几何规律的利用仍不充分；其二，在线定位对先验地图结构信息和不同方向观测差异的利用仍有提升空间；其三，大尺度先验地图的组织与动态调用尚缺乏面向矿井连续定位任务的系统化研究。上述问题共同构成了本文的研究切入点。

\section{研究内容和结构安排}

针对矿井巷道场景中高精地图构建、先验约束定位以及大尺度地图调用三方面存在的问题，本文围绕“多源结构先验如何在离线建图、在线定位和地图服务中持续发挥作用”这一主线展开研究，整体研究内容包括以下三个方面。

第一，面向矿井巷道离线建图问题，研究基于几何结构约束的位姿图建图方法。针对长距离运行过程中误差累积明显、局部观测约束不足的问题，本文在关键帧位姿图框架下引入巷道长度关系、方向关系以及局部平面辅助信息，构建多因子联合优化模型，以提升先验地图的全局一致性和结构准确性。

第二，面向高精先验地图条件下的在线定位问题，研究基于先验体素约束与轴向分解融合的定位方法。针对矿井巷道场景中轴向与非轴向约束不均衡、地图观测可靠性波动以及结构重复区域定位不稳定等问题，本文构建帧间配准与地图配准的双源观测模型，并结合先验体素命中统计和巷道中轴信息，对不同方向上的观测实施差异化融合，以提高定位精度与鲁棒性。

第三，面向大尺度先验地图服务问题，研究先验地图组织与动态加载方法。针对整图加载开销大、长期运行资源压力高的问题，本文对离线高精地图进行规则瓦片化组织，建立元数据索引，并依据当前位置实施局部地图按需生成与动态调用，从而提高先验地图在连续定位任务中的使用效率。

围绕上述研究内容，本文形成了从“高精地图构建”到“先验约束定位”再到“地图组织调用”的技术路线。其基本思路是：首先通过离线建图获得具有较好全局一致性的高精先验地图；随后在在线定位阶段利用地图中的体素占据关系和巷道结构方向信息增强定位稳定性；最后将先验地图进一步组织为适于在线调用的结构化形式，为大尺度连续定位提供地图服务支撑。该技术路线使得多源结构先验不再局限于单一环节，而是在建图、定位与系统运行层面形成连续作用链条。

在上述研究内容基础上，全文结构安排如下。

全文共分为六章，各章内容安排如下。

第 1 章为绪论，阐述论文的研究背景与意义，综述矿井巷道建图、先验地图定位以及大尺度地图组织与调用等方面的研究现状，明确本文的研究内容、技术路线和章节安排。

第 2 章为系统总体框架与相关理论基础，介绍本文研究所依托的系统框架，并对三维空间运动描述、点云配准、状态估计、图优化与因子图等理论基础进行统一说明，为后续章节的方法建模提供理论支撑。

第 3 章研究基于几何结构约束的矿井巷道建图方法。在关键帧位姿图框架下，建立全局长度约束、全局角度约束以及局部平面辅助约束等多类几何因子，通过多因子联合优化提高矿井巷道先验地图的全局一致性和结构准确性。

第 4 章研究基于先验体素约束与轴向分解融合的矿井定位方法。围绕位姿预测、帧间配准、地图配准和差异化融合估计等关键环节，构建适用于矿井巷道场景的先验约束定位方法，并通过真实地下矿井数据集和 Gazebo 仿真数据集进行实验验证。

第 5 章研究面向连续定位任务的高精先验地图组织与动态加载方法。在统一先验地图和定位方法条件下，对先验地图进行规则网格划分和元数据索引构建，建立局部地图动态调用机制，并分析整图加载与动态加载两种方式在系统运行性能上的差异。

第 6 章对全文工作进行总结，概括本文的主要研究成果、特点与不足，并对后续研究方向进行展望。
